% ============================================================================
% APPENDIX OT: LIT-OTHER - Miscellaneous Research Integration
% ============================================================================
% Category: LIT (Literature Integration)
% Version: 1.0
% Last Updated: 2026-01-20
% Dependencies: All LIT appendices
% Language: english
% ============================================================================

% ----------------------------------------------------------------------------
% HEADER BLOCK
% ----------------------------------------------------------------------------
% Code: OT
% Category: LIT (Literature Integration)
% Title: Miscellaneous Research Integration
% Author: EBF Framework
% Status: Active (Growing Collection)
% Priority: Medium
% Evidence-Base: Variable (see data/paper-sources.yaml)
% ----------------------------------------------------------------------------

\chapter{LIT-OTHER: Miscellaneous Research Integration}
\label{app:lit-other}


\begin{tcolorbox}[colback=blue!5!white, colframe=blue!75!black,
    title=Quick Reference: Appendix OT]

\textbf{Appendix:} OT (LIT)

\textbf{Core Question:} What are the key findings in this domain?

\textbf{Key Concepts:}
\begin{itemize}[nosep]
\item Framework integration
\item Parameter validation
\item Cross-references
\end{itemize}

\textbf{Cross-References:} See related appendices below
\end{tcolorbox}

% ============================================================================
% CROSS-REFERENCE MAP
% ============================================================================
\begin{tcolorbox}[colback=blue!5!white,colframe=blue!75!black,title=Cross-Reference Map]
\textbf{Outgoing References:}
\begin{itemize}
    \item All CORE appendices: Thematic connections vary by paper
    \item All LIT appendices: Complementary research
    \item All DOMAIN appendices: Application areas
\end{itemize}

\textbf{Incoming References:}
\begin{itemize}
    \item New papers added as EBF-relevant research emerges
    \item Papers may migrate to dedicated LIT appendices when author corpus grows
\end{itemize}
\end{tcolorbox}

% ============================================================================
% CHAPTER LINKAGE
% ============================================================================
\begin{tcolorbox}[colback=green!5!white,colframe=green!75!black,title=Chapter Linkage]
\textbf{Primary Chapter:} Variable (depends on paper topic)

\textbf{Purpose:} Collection point for EBF-relevant papers that do not yet have a dedicated author-based LIT appendix.
\end{tcolorbox}

% ============================================================================
% ABSTRACT AND QUICK REFERENCE
% ============================================================================
\begin{abstract}
This appendix serves as a collection point for important research contributions that are relevant to the EBF framework but do not fit into existing author-based LIT appendices. Papers collected here may eventually migrate to dedicated appendices as author corpora grow or new thematic clusters emerge. The appendix is organized thematically rather than by author, allowing flexible integration of emerging research, cross-disciplinary contributions, and recent working papers.
\end{abstract}

\begin{tcolorbox}[colback=gray!5!white,colframe=gray!75!black,title=Quick Reference]
\textbf{Collection Criteria:}
\begin{itemize}
    \item Directly relevant to EBF framework
    \item No existing dedicated author appendix
    \item High-quality peer-reviewed or NBER working papers
    \item Fills gap in current LIT coverage
\end{itemize}

\textbf{Current Thematic Clusters:}
\begin{itemize}
    \item Market morality and repugnant transactions
    \item Choice architecture and decision complexity
    \item Cognitive psychology foundations
    \item Prosocial motivation and work meaning
    \item Organizational psychology
    \item Environmental behavior \& sustainable consumption
    \item Rational inattention \& behavioral inattention (NTM)
    \item Salience theory (NTM)
    \item Procrastination \& time inconsistency (NTM)
    \item Habit formation \& automaticity (NTM)
    \item Intent-action gap (NTM)
    \item Social norms \& norm change (NTM)
    \item Political psychology \& authoritarianism (NEW)
    \item Policy scaling \& implementation
    \item Network science \& social contagion (UNTCM)
    \item Complexity \& path dependence (UNTCM)
    \item Econometrics \& causal inference (UNTCM)
    \item Tipping points \& critical transitions (UNTCM)
    \item Learning \& expectations (UNTCM)
    \item Behavioral science methodology (UNTCM)
\end{itemize}

\textbf{Migration Policy:} When 10+ papers from a single author or coherent research program accumulate, consider creating a dedicated LIT appendix.
\end{tcolorbox}

% ============================================================================
% FUNDAMENTAL QUESTION
% ============================================================================
% -----------------------------------------------------------------------------
% APPENDIX SCOPE BOX
% -----------------------------------------------------------------------------

\begin{tcolorbox}[
    colback=orange!5!white,
    colframe=orange!75!black,
    title={\textbf{Appendix Scope: Ziel / In-Scope / Out-of-Scope / Constraints}},
    fonttitle=\bfseries
]

\textbf{Ziel (Objective):}
\begin{itemize}[nosep]
    \item What emerging research contributions extend or challenge the EBF framework, and how should they be integrated before dedicated appendices are warranted?
\end{itemize}

\textbf{In-Scope:}
\begin{itemize}[nosep]
    \item Fundamental Question
    \item Thematic Clusters
    \item Integration Guidelines
\end{itemize}

\textbf{Out-of-Scope (delegiert an andere Appendices):}
\begin{itemize}[nosep]
    \item REF-GLOSSARY $\rightarrow$ Appendix G
    \item REF-GLOSSARY $\rightarrow$ Appendix G
\end{itemize}

\textbf{Constraints (Anwendungsgrenzen):}
\begin{itemize}[nosep]
    \item [TODO: Define when this appendix does NOT apply]
    \item [TODO: Define required prerequisites]
    \item [TODO: Define known limitations]
\end{itemize}

\textbf{Lieferobjekte (Deliverables):}
\begin{itemize}[nosep]
    \item [TODO: Add deliverables]
\end{itemize}

\end{tcolorbox}


\section{Fundamental Question}
\label{sec:ot-fundamental}

\begin{tcolorbox}[colback=red!5!white,colframe=red!75!black,title=Fundamental Question]
\textbf{What emerging research contributions extend or challenge the EBF framework, and how should they be integrated before dedicated appendices are warranted?}
\end{tcolorbox}

% ============================================================================
% THEMATIC CLUSTERS
% ============================================================================
\section{Thematic Clusters}
\label{sec:ot-clusters}

\subsection{Market Morality and Repugnant Transactions}
\label{subsec:ot-repugnant}

Research on contested markets where moral considerations constrain or shape economic exchange.

\begin{tcolorbox}[colback=yellow!5!white,colframe=yellow!75!black,title=Key Paper: Elias, Lacetera \& Macis (2026)]
\textbf{Title:} The Morality of Market Exchanges: Between Societal Values and Tradeoffs

\textbf{Source:} NBER Working Paper 34647

\textbf{Key Contributions:}
\begin{itemize}
    \item Framework for analyzing contested/repugnant markets
    \item Two case studies: (1) compensated organ donation and surrogacy, (2) price gouging in emergencies
    \item Tradeoffs between autonomy, fairness, and efficiency
    \item Institutional features that separate compensation from allocation
\end{itemize}

\textbf{EBF Connections:}
\begin{itemize}
    \item Bénabou-Tirole (RB/JT): Taboos, sacred values, motivation crowding
    \item Fehr (K): Fairness norms in markets
    \item Stratification Economics (SE): Distributional concerns
\end{itemize}

\textbf{Evidence Tier:} 2-3 (survey evidence, case studies)
\end{tcolorbox}

\textbf{Related Research (potential additions):}
\begin{itemize}
    \item Roth, A. E. (2007). Repugnance as a Constraint on Markets. \textit{JEP}.
    \item Satz, D. (2010). \textit{Why Some Things Should Not Be For Sale}. Oxford.
    \item Sandel, M. (2012). \textit{What Money Can't Buy}. FSG.
    \item Ambuehl, S., Niederle, M., \& Roth, A. E. (2015). More Money, More Problems?
\end{itemize}

\subsection{Choice Architecture and Decision Complexity}
\label{subsec:ot-choice-architecture}

Foundational research on how choice set design affects decision quality and satisfaction.

\begin{tcolorbox}[colback=yellow!5!white,colframe=yellow!75!black,title=Key Paper: Iyengar \& Lepper (2000)]
\textbf{Title:} When Choice is Demotivating: Can One Desire Too Much of a Good Thing?

\textbf{Source:} Journal of Personality and Social Psychology, 79(6), 995--1006.

\textbf{Key Contributions (``Jam Study''):}
\begin{itemize}
    \item Field experiment: 24 jam varieties vs. 6 jam varieties
    \item Extensive choice (24): 3\% purchase rate
    \item Limited choice (6): 30\% purchase rate
    \item Established ``choice overload'' as empirical phenomenon
\end{itemize}

\textbf{EBF Connections:}
\begin{itemize}
    \item Thaler/Sunstein (THL/SUN): Choice architecture, default effects
    \item METHOD-TOOLKIT (HHH): Menu sizing guidelines ($n^* \approx 7$)
    \item DOMAIN-HR: Benefits package design (INT-HR-006)
\end{itemize}

\textbf{Evidence Tier:} 1 (field experiment, highly replicated)

\textbf{Key Parameter:} Choice overload threshold: $n^* = 6$--7 options
\end{tcolorbox}

\begin{tcolorbox}[colback=yellow!5!white,colframe=yellow!75!black,title=Key Paper: Schwartz (2004)]
\textbf{Title:} The Paradox of Choice: Why More Is Less

\textbf{Source:} Harper Collins (Book).

\textbf{Key Contributions:}
\begin{itemize}
    \item Synthesizes research on choice overload effects
    \item Identifies four mechanisms: (1) Choice paralysis, (2) Opportunity cost salience, (3) Escalation of expectations, (4) Self-blame for suboptimal choices
    \item Proposes satisfaction decay formula: $Z(n) = Z_{max} \times (1 - \beta \times \max(0, n - n^*))$
\end{itemize}

\textbf{EBF Connections:}
\begin{itemize}
    \item Extends Iyengar \& Lepper to theoretical framework
    \item Foundation for benefits choice architecture
    \item Connects to Kahneman's System 1/System 2 (cognitive load)
\end{itemize}

\textbf{Evidence Tier:} 2 (synthesis, popular science)

\textbf{Key Parameter:} Satisfaction decay: $\beta \approx 0.03$ per option above threshold
\end{tcolorbox}

\subsection{Cognitive Psychology Foundations}
\label{subsec:ot-cognitive}

Classic contributions from cognitive psychology that underpin behavioral economics.

\begin{tcolorbox}[colback=yellow!5!white,colframe=yellow!75!black,title=Key Paper: Miller (1956)]
\textbf{Title:} The Magical Number Seven, Plus or Minus Two: Some Limits on Our Capacity for Processing Information

\textbf{Source:} Psychological Review, 63(2), 81--97.

\textbf{Key Contributions:}
\begin{itemize}
    \item Establishes working memory capacity: $7 \pm 2$ items
    \item Foundation for ``chunking'' as cognitive strategy
    \item Explains limits on simultaneous option processing
    \item Most cited paper in cognitive psychology
\end{itemize}

\textbf{EBF Connections:}
\begin{itemize}
    \item Theoretical foundation for choice overload (Iyengar, Schwartz)
    \item Constrains optimal menu size in choice architecture
    \item Links to CORE-HOW bounded rationality
    \item Gigerenzer (GI): Ecological rationality, heuristics
\end{itemize}

\textbf{Evidence Tier:} 1 (foundational, universally replicated)

\textbf{Key Parameter:} Working memory capacity: $n = 7 \pm 2$
\end{tcolorbox}

\subsection{Prosocial Motivation and Work Meaning}
\label{subsec:ot-prosocial}

Research on how purpose and prosocial impact affect motivation and performance.

\begin{tcolorbox}[colback=yellow!5!white,colframe=yellow!75!black,title=Key Paper: Grant (2008)]
\textbf{Title:} Does Intrinsic Motivation Fuel the Prosocial Fire? Motivational Synergy in Predicting Persistence, Performance, and Productivity

\textbf{Source:} Journal of Applied Psychology, 93(1), 48--58.

\textbf{Key Contributions:}
\begin{itemize}
    \item Intrinsic and prosocial motivation interact multiplicatively
    \item Neither alone sufficient for sustained performance
    \item Field studies: call center fundraisers, firefighters
    \item Prosocial identity increases persistence under difficulty
\end{itemize}

\textbf{EBF Connections:}
\begin{itemize}
    \item Extends Deci/Ryan self-determination theory
    \item Connects to Akerlof \& Kranton identity economics (GA)
    \item Foundation for T5 (Identity) interventions
    \item Supports Mental Identity Budgeting (INT-HR-007)
\end{itemize}

\textbf{Evidence Tier:} 1 (field studies, longitudinal)

\textbf{Key Parameter:} Prosocial identity multiplier: $\sigma = 1.5$--1.8
\end{tcolorbox}

\subsection{Organizational Psychology}
\label{subsec:ot-organizational}

Contributions from organizational psychology on employee commitment and retention.

\begin{tcolorbox}[colback=yellow!5!white,colframe=yellow!75!black,title=Key Paper: Meyer \& Allen (1991)]
\textbf{Title:} A Three-Component Conceptualization of Organizational Commitment

\textbf{Source:} Human Resource Management Review, 1(1), 61--89.

\textbf{Key Contributions:}
\begin{itemize}
    \item Three types of organizational commitment:
    \begin{enumerate}
        \item \textbf{Affective:} Emotional attachment (identity-based)
        \item \textbf{Continuance:} Cost of leaving (switching costs)
        \item \textbf{Normative:} Obligation to stay (social norms)
    \end{enumerate}
    \item Validated scales used in 1000+ studies
    \item Predicts turnover, absenteeism, performance
\end{itemize}

\textbf{EBF Connections:}
\begin{itemize}
    \item Affective commitment $\leftrightarrow$ T5 (Identity) interventions
    \item Continuance commitment $\leftrightarrow$ T7 (Financial) + switching costs
    \item Normative commitment $\leftrightarrow$ T6 (Social) norms
    \item Maps to HR Retention Portfolio (PORT-HR-001)
\end{itemize}

\textbf{Evidence Tier:} 1 (foundational, extensively validated)

\textbf{Key Insight:} Identity Lock-In activates affective + continuance commitment simultaneously
\end{tcolorbox}

\subsection{Emerging Behavioral Economics}
\label{subsec:ot-emerging}

Recent contributions extending core behavioral economics insights.

\textit{[Papers to be added as identified]}

\subsection{Rational Inattention \& Behavioral Inattention}
\label{subsec:ot-inattention}

\textit{Added 2026-01-19: Supporting NTM Awareness Decay Mechanisms}

\begin{tcolorbox}[colback=yellow!5!white,colframe=yellow!75!black,title=Key Paper: Sims (2003)]
\textbf{Citation:} Sims, Christopher A. (2003). ``Implications of Rational Inattention.'' \textit{Journal of Monetary Economics}, 50(3), 665--690.

\textbf{Core Contribution:}
\begin{itemize}[nosep]
    \item Rational inattention: attention is scarce cognitive resource
    \item Information processing capacity constraints drive behavior
    \item Explains sticky prices, delayed responses, information neglect
\end{itemize}

\textbf{NTM Relevance:} Theoretical foundation for awareness decay (NTM-A1): limited attention capacity means awareness naturally diminishes without reinforcement.

\textbf{10C Integration:}
\begin{itemize}[nosep]
    \item \textbf{Domain:} Information Economics, Macroeconomics
    \item \textbf{Dimension:} E (Epistemic), P (Physical)
    \item \textbf{CORE:} AWARE (AU) - attention constraints
    \item \textbf{Complementarity:} γ = 0.65 (attention-action coupling)
\end{itemize}
\end{tcolorbox}

\subsection{Salience Theory}
\label{subsec:ot-salience}

\textit{Added 2026-01-19: Supporting NTM Intent Salience Decay}

\begin{tcolorbox}[colback=yellow!5!white,colframe=yellow!75!black,title=Key Paper: Bordalo et al. (2012)]
\textbf{Citation:} Bordalo, Pedro, Gennaioli, Nicola, and Shleifer, Andrei (2012). ``Salience Theory of Choice Under Risk.'' \textit{QJE}, 127(3), 1243--1285.

\textbf{Core Contribution:}
\begin{itemize}[nosep]
    \item Context-dependent salience: extreme payoffs attract attention
    \item Local thinking: decisions based on salient subset of information
    \item Explains framing effects, risk preferences, consumer choice anomalies
\end{itemize}

\textbf{NTM Relevance:} Foundation for intent salience decay $(1-\eta)^n$ in NTM-T1: without salient triggers, behavioral intentions fade from awareness.

\textbf{10C Integration:}
\begin{itemize}[nosep]
    \item \textbf{Domain:} Decision Theory, Risk
    \item \textbf{Dimension:} E (Epistemic), F (Financial)
    \item \textbf{CORE:} AWARE (AU) - salience-driven attention
    \item \textbf{Complementarity:} γ = 0.72 (salience-choice coupling)
\end{itemize}
\end{tcolorbox}

\subsection{Procrastination \& Time Inconsistency}
\label{subsec:ot-procrastination}

\textit{Added 2026-01-19: Supporting NTM Threshold Drift}

\begin{tcolorbox}[colback=yellow!5!white,colframe=yellow!75!black,title=Key Paper: O'Donoghue \& Rabin (2001)]
\textbf{Citation:} O'Donoghue, Ted and Rabin, Matthew (2001). ``Choice and Procrastination.'' \textit{QJE}, 116(1), 121--160.

\textbf{Core Contribution:}
\begin{itemize}[nosep]
    \item Naive vs. sophisticated present-biased agents
    \item Procrastination as rational response to time inconsistency
    \item Deadlines and commitment devices as solutions
\end{itemize}

\textbf{NTM Relevance:} Foundation for threshold drift $\theta(t) = \theta_0 + \delta \cdot t$ in NTM-T1: procrastination accumulates over time, raising action barriers.

\textbf{10C Integration:}
\begin{itemize}[nosep]
    \item \textbf{Domain:} Time Preferences, Self-Control
    \item \textbf{Dimension:} P (Physical behavior)
    \item \textbf{CORE:} WHEN (V) - temporal dynamics, READY (AV) - willingness threshold
    \item \textbf{Complementarity:} γ = 0.58 (time-action coupling)
\end{itemize}
\end{tcolorbox}

\subsection{Habit Formation \& Automaticity}
\label{subsec:ot-habits}

\textit{Added 2026-01-19: Supporting NTM Long-Term Lock-In}

\begin{tcolorbox}[colback=yellow!5!white,colframe=yellow!75!black,title=Key Paper: Wood \& Rünger (2016)]
\textbf{Citation:} Wood, Wendy and Rünger, Dennis (2016). ``Psychology of Habit.'' \textit{Annual Review of Psychology}, 67, 289--314.

\textbf{Core Contribution:}
\begin{itemize}[nosep]
    \item Habit formation requires context stability + repetition
    \item Automaticity reduces cognitive load but creates lock-in
    \item Breaking habits requires environmental disruption
\end{itemize}

\textbf{NTM Relevance:} Foundation for long-term equilibrium lock-in in NTM-A1: habitual behaviors become automatic, resistant to change.

\textbf{10C Integration:}
\begin{itemize}[nosep]
    \item \textbf{Domain:} Psychology, Behavior Change
    \item \textbf{Dimension:} P (Physical), D (Development)
    \item \textbf{CORE:} STAGE (AW) - maintenance phase, READY (AV) - automaticity
    \item \textbf{Complementarity:} γ = 0.80 (context-habit coupling)
\end{itemize}
\end{tcolorbox}

\begin{tcolorbox}[colback=gray!5!white,colframe=gray!75!black,title=Supporting Paper: Verplanken \& Orbell (2003)]
\textbf{Citation:} Verplanken, Bas and Orbell, Sheila (2003). ``Reflections on Past Behavior: A Self-Report Index of Habit Strength.'' \textit{Journal of Applied Social Psychology}, 33(6), 1313--1330.

\textbf{Core Contribution:} Self-Report Habit Index (SRHI) - validated measure of habit strength based on automaticity, frequency, and self-identity.

\textbf{NTM Relevance:} Operationalization of habit lock-in measurement.
\end{tcolorbox}

\subsection{Intent-Action Gap}
\label{subsec:ot-intent-action}

\textit{Added 2026-01-19: Supporting NTM Temporal Dynamics}

\begin{tcolorbox}[colback=yellow!5!white,colframe=yellow!75!black,title=Key Paper: Sheeran \& Webb (2016)]
\textbf{Citation:} Sheeran, Paschal and Webb, Thomas L. (2016). ``The Intention-Behavior Gap.'' \textit{Social and Personality Psychology Compass}, 10(9), 503--518.

\textbf{Core Contribution:}
\begin{itemize}[nosep]
    \item Medium-large gap between intention and behavior (d = 0.36 meta-analytic)
    \item Gap widens over time without implementation planning
    \item Intention stability, behavioral control, and implementation intentions as moderators
\end{itemize}

\textbf{NTM Relevance:} Foundation for willingness decay and threshold drift: even strong intentions decay without supporting structures.

\textbf{10C Integration:}
\begin{itemize}[nosep]
    \item \textbf{Domain:} Health Psychology, Behavior Change
    \item \textbf{Dimension:} P (Physical), D (Development)
    \item \textbf{CORE:} READY (AV) - willingness, STAGE (AW) - preparation to action gap
    \item \textbf{Complementarity:} γ = 0.55 (intent-behavior coupling)
\end{itemize}
\end{tcolorbox}

\subsection{Social Norms \& Norm Change}
\label{subsec:ot-norms}

\textit{Added 2026-01-19: Supporting NTM Context Drift}

\begin{tcolorbox}[colback=yellow!5!white,colframe=yellow!75!black,title=Key Paper: Bicchieri (2005)]
\textbf{Citation:} Bicchieri, Cristina (2005). ``The Grammar of Society: The Nature and Dynamics of Social Norms.'' \textit{Cambridge University Press}.

\textbf{Core Contribution:}
\begin{itemize}[nosep]
    \item Social norms = conditional preferences + empirical/normative expectations
    \item Norm change requires changing expectations, not just preferences
    \item Tipping point dynamics: small changes can cascade
\end{itemize}

\textbf{NTM Relevance:} Foundation for context drift $\Psi(t + \Delta t) = \Psi(t) + \mu_{\Psi} \cdot \Delta t$ in NTM-A2: social norms evolve autonomously, creating positive or negative context drift.

\textbf{10C Integration:}
\begin{itemize}[nosep]
    \item \textbf{Domain:} Social Philosophy, Game Theory
    \item \textbf{Dimension:} S (Social)
    \item \textbf{CORE:} WHEN (V) - social context, WHO (AAA) - collective preferences
    \item \textbf{Complementarity:} γ = 0.75 (norm-behavior coupling)
\end{itemize}
\end{tcolorbox}

\subsection{Environmental Behavior \& Sustainable Consumption}
\label{subsec:ot-environmental}

Research on sustainable consumer behavior, environmental attitudes, and the attitude-behavior gap.
\textit{Added 2026-01-31: Supporting CORE-AWARE and cognitive dissonance in consumption}

\begin{tcolorbox}[colback=yellow!5!white,colframe=yellow!75!black,title=Key Paper: Giner, Nauges \& Hassett (2025)]
\textbf{Citation:} Giner, C\'eline, Nauges, C\'eline, and Hassett, Katherine (2025). ``Patterns in sustainable food choices and policy support: Novel evidence from nine countries.'' \textit{Food Policy}, 128, 102748.

\textbf{Core Contribution:}
\begin{itemize}[nosep]
    \item Latent Class Analysis identifies 4 consumer segments across 9 OECD countries (n=8,261)
    \item Documents ``Meat Paradox'' (Class 3, 26\%): High environmental concern but high meat consumption
    \item Policy support varies dramatically by class: meat tax support ranges from 16\% (Class 1) to 36\% (Class 4)
    \item Cross-country analysis reveals universal vs. culture-specific patterns
\end{itemize}

\textbf{EBF Connections:}
\begin{itemize}[nosep]
    \item \textbf{Cognitive Dissonance (MS-NE-007):} Class 3 exemplifies attitude-behavior gap
    \item \textbf{CORE-AWARE (AU):} Awareness of environmental impact $\neq$ behavior change
    \item \textbf{CORE-WHEN (V):} Cultural context modulates policy acceptance
    \item \textbf{Social Norms (MS-SN-001):} Peer effects in consumption patterns
\end{itemize}

\textbf{Key Parameters:}
\begin{itemize}[nosep]
    \item $\pi_{class}$: Consumer segment distribution [0.41, 0.21, 0.26, 0.12]
    \item CDI = 0.26: Cognitive Dissonance Index (share showing attitude-behavior gap)
    \item $\alpha_{tax}$: Meat tax support [0.16, 0.24, 0.28, 0.36] by class
    \item $\rho_{concern}$: Environmental concern-behavior correlation gap (0.06)
\end{itemize}

\textbf{10C Integration:}
\begin{itemize}[nosep]
    \item \textbf{Domain:} Consumer Behavior, Environmental Policy
    \item \textbf{Dimension:} E (Environmental), P (Physical consumption), D (Development)
    \item \textbf{CORE:} AWARE (AU) - awareness-action gap, WHEN (V) - cultural context
    \item \textbf{Complementarity:} $\gamma = 0.06$ (concern-behavior coupling, weak)
\end{itemize}

\textbf{Evidence Tier:} 2 (cross-country survey, OECD EPIC, peer-reviewed)
\end{tcolorbox}

\subsection{Political Psychology \& Authoritarianism}
\label{subsec:ot-political-psychology}

Research on right-wing authoritarianism (RWA), social dominance orientation (SDO), conspiracy beliefs, and their behavioral manifestations.
\textit{Added 2026-02-02: Supporting CAT-21 Political Psychology theories (MS-PP-001 to MS-PP-005)}

\begin{tcolorbox}[colback=yellow!5!white,colframe=yellow!75!black,title=Key Paper: DÖW Rechtsextremismus-Barometer (2024)]
\textbf{Citation:} \citet{PAP-doew2024rechtsextremismus}. ``Rechtsextremismus-Barometer Österreich 2024.'' Dokumentationsarchiv des österreichischen Widerstandes, Vienna.

\textbf{Core Contribution:}
\begin{itemize}[nosep]
    \item First comprehensive Austrian extremism survey (n=2,198, representative)
    \item 10\% of population classified as extreme right (high RWA + SDO)
    \item 47\% believe in ``Great Replacement'' (highest in Europe)
    \item 50\% support comprehensive remigration policies
    \item 58\% of extreme right segment votes FPÖ
\end{itemize}

\textbf{EBF Connections:}
\begin{itemize}[nosep]
    \item \textbf{RWA (MS-PP-001):} Altemeyer's Right-Wing Authoritarianism scale
    \item \textbf{SDO (MS-PP-002):} Sidanius-Pratto Social Dominance Orientation
    \item \textbf{Conspiracy Mentality (MS-PP-003):} Generalized conspiracy belief tendency
    \item \textbf{CORE-WHO (AAA):} Defines hard-core FPÖ segment (~6\% of total population)
    \item \textbf{CORE-WHAT (C):} Identity-driven utility (U\_IDN dominant over U\_F)
\end{itemize}

\textbf{Key Parameters:}
\begin{itemize}[nosep]
    \item $\pi_{extreme} = 0.10$: Extreme right prevalence (PAR-PP-001)
    \item $\pi_{replacement} = 0.47$: Great Replacement belief (PAR-PP-002)
    \item $\pi_{remigration} = 0.50$: Remigration support (PAR-PP-003)
    \item $\rho_{fpoe} = 0.58$: FPÖ share among extreme right (PAR-PP-004)
\end{itemize}

\textbf{10C Integration:}
\begin{itemize}[nosep]
    \item \textbf{Domain:} Political Psychology, Austrian Politics
    \item \textbf{Dimension:} E (Identity/Ego), S (In-group preference)
    \item \textbf{CORE:} WHO (AAA) - segment definition, WHAT (C) - identity utility
    \item \textbf{Complementarity:} $\gamma = 0.75$ (RWA $\times$ SDO $\times$ Conspiracy reinforcing)
\end{itemize}

\textbf{Evidence Tier:} 2 (representative survey, DÖW institutional credibility)

\textbf{Case:} CAS-890 (Rechtsextremismus-Segment Österreich 2024)
\end{tcolorbox}

\begin{tcolorbox}[colback=gray!10!white,colframe=gray!75!black,title=Related Papers in This Cluster]
\begin{itemize}[nosep]
    \item \citet{PAP-kleppestoe2024genetic}: Genetic basis of RWA/SDO (30-50\% heritable)
    \item \citet{PAP-herold2025replacement}: Cross-European Great Replacement beliefs
    \item \citet{PAP-ifes2022antisemitismus}: Austrian antisemitism survey (31\% prevalence)
    \item \citet{PAP-weisskircher2025normalised}: FPÖ normalization trajectory
\end{itemize}
\end{tcolorbox}

\subsection{Policy Scaling \& Implementation}
\label{subsec:ot-policy-scaling}

Research on scaling behavioral interventions from pilot to national/global level.
\textit{Added 2026-01-19: Supporting Portfolio Archetypes \& Sustainable Implementation}

\begin{tcolorbox}[colback=yellow!5!white,colframe=yellow!75!black,title=Key Paper: Soman \& Hossain (2020)]
\textbf{Citation:} Soman, Dilip and Hossain, Tamim (2020). ``Successfully Scaling Nudges at the National Level: A Framework and Case Study.'' \textit{Behavioural Public Policy}, Cambridge University Press, 1--25.

\textbf{Core Contribution:}
\begin{itemize}[nosep]
    \item Pilot-to-scale decay factor: 0.3--0.5 (effect sizes diminish at scale)
    \item Cultural adaptation required for cross-context implementation
    \item Framework for national-level nudge deployment
\end{itemize}

\textbf{10C Integration:}
\begin{itemize}[nosep]
    \item \textbf{Domain:} Policy, Government
    \item \textbf{Dimension:} P (Physical behavior), S (Social scaling)
    \item \textbf{Context:} National programs, developing country (India)
    \item \textbf{Complementarity:} γ = 0.50 (context-adaptation synergy)
\end{itemize}

\textbf{EBF Use:}
\begin{itemize}[nosep]
    \item Portfolio Archetype ``Sustainable'' (Chapter 20, Section 4.4)
    \item FEPSDE-KPI scaling adjustments (Chapter 20, Section 4.5)
    \item HHH 3×3 KPI Protocol (Definition HHH-D7)
\end{itemize}
\end{tcolorbox}

% ============================================================================
% UNTCM EVIDENCE CLUSTERS (Added 2026-01-20)
% Supporting Unified NTM-Context Model (Appendix PBB)
% ============================================================================

\subsection{Network Science \& Social Contagion}
\label{subsec:ot-network-science}

Research on how behavior spreads through social networks and the role of network topology.
\textit{Added 2026-01-20: Supporting UNTCM-X1 (Network-Dependent Feedback)}

\begin{tcolorbox}[colback=yellow!5!white,colframe=yellow!75!black,title=Key Paper: Centola (2018)]
\textbf{Citation:} Centola, Damon (2018). ``How Behavior Spreads: The Science of Complex Contagions.'' Princeton University Press.

\textbf{Core Contribution:}
\begin{itemize}[nosep]
    \item Complex contagions require social reinforcement (multiple exposures)
    \item Clustered networks outperform random networks for complex behaviors by 2-4x
    \item Distinction between simple contagion (information) and complex contagion (behavior)
\end{itemize}

\textbf{UNTCM Integration:}
\begin{itemize}[nosep]
    \item X1: $\kappa_{eff} = \kappa_0 \cdot C_G^{0.5}$ for complex behaviors
    \item Seeding strategy: Target clustered subgroups for T5/T6 interventions
\end{itemize}
\end{tcolorbox}

\begin{tcolorbox}[colback=yellow!5!white,colframe=yellow!75!black,title=Key Paper: Centola \& Macy (2007)]
\textbf{Citation:} Centola, Damon and Macy, Michael (2007). ``Complex Contagions and the Weakness of Long Ties.'' \textit{American Journal of Sociology}, 113(3), 702--734.

\textbf{Core Contribution:}
\begin{itemize}[nosep]
    \item Granovetter's ``strength of weak ties'' applies only to simple contagions
    \item Long ties can fragment adoption of complex behaviors
    \item Local clustering essential for behavioral reinforcement
\end{itemize}
\end{tcolorbox}

\begin{tcolorbox}[colback=yellow!5!white,colframe=yellow!75!black,title=Key Paper: Watts (2002)]
\textbf{Citation:} Watts, Duncan J (2002). ``A Simple Model of Global Cascades on Random Networks.'' \textit{PNAS}, 99(9), 5766--5771.

\textbf{Core Contribution:}
\begin{itemize}[nosep]
    \item Cascade window: 10-25\% initial adoption can trigger global cascade
    \item Network degree distribution affects cascade threshold
    \item Phase transition between no-cascade and full-cascade regimes
\end{itemize}

\textbf{UNTCM Integration:}
\begin{itemize}[nosep]
    \item X1: Tipping point $\bar{B}^* \in [0.10, 0.25]$ depending on topology
    \item E2: Cascade probability as function of network parameters
\end{itemize}
\end{tcolorbox}

\begin{tcolorbox}[colback=orange!5!white,colframe=orange!75!black,title=CONTRA Paper: Ugander et al. (2012)]
\textbf{Citation:} Ugander, Johan et al. (2012). ``Structural Diversity in Social Contagion.'' \textit{PNAS}, 109(16), 5962--5966.

\textbf{Core Contribution (CONTRA-X1):}
\begin{itemize}[nosep]
    \item Structural diversity (variety of connections) matters more than quantity
    \item Simple topology metrics may oversimplify real network effects
    \item Facebook data shows nuanced adoption dynamics
\end{itemize}

\textbf{Implication:} X1 may need diversity index extension beyond clustering coefficient.
\end{tcolorbox}

\subsection{Complexity \& Path Dependence}
\label{subsec:ot-complexity}

Research on increasing returns, lock-in, and path-dependent dynamics.
\textit{Added 2026-01-20: Supporting UNTCM-X2 (Hysteresis)}

\begin{tcolorbox}[colback=yellow!5!white,colframe=yellow!75!black,title=Key Paper: Arthur (1989)]
\textbf{Citation:} Arthur, W Brian (1989). ``Competing Technologies, Increasing Returns, and Lock-In by Historical Events.'' \textit{Economic Journal}, 99(394), 116--131.

\textbf{Core Contribution:}
\begin{itemize}[nosep]
    \item Small historical events can determine long-run outcomes
    \item Increasing returns create multiple equilibria
    \item Lock-in occurs even without inherent superiority (QWERTY example)
\end{itemize}

\textbf{UNTCM Integration:}
\begin{itemize}[nosep]
    \item X2: $\bar{B}^*_{down} \ll \bar{B}^*_{up}$ (asymmetric tipping)
    \item Once in Regime II, system is robust to negative shocks
\end{itemize}
\end{tcolorbox}

\begin{tcolorbox}[colback=orange!5!white,colframe=orange!75!black,title=CONTRA Paper: Liebowitz \& Margolis (1995)]
\textbf{Citation:} Liebowitz, Stan J and Margolis, Stephen E (1995). ``Path Dependence, Lock-In, and History.'' \textit{Journal of Law, Economics, and Organization}, 11(1), 205--226.

\textbf{Core Contribution (CONTRA-X2):}
\begin{itemize}[nosep]
    \item Path dependence may be overstated
    \item Markets often correct ``inefficient'' lock-ins
    \item QWERTY example questioned empirically
\end{itemize}

\textbf{Implication:} X2 hysteresis may be weaker in competitive markets; applies more to identity/habit domains.
\end{tcolorbox}

\subsection{Econometrics \& Causal Inference}
\label{subsec:ot-econometrics}

Research on regime-switching models, causal inference, and experimental design.
\textit{Added 2026-01-20: Supporting UNTCM-X4, E1, E3}

\begin{tcolorbox}[colback=yellow!5!white,colframe=yellow!75!black,title=Key Paper: Hamilton (1989)]
\textbf{Citation:} Hamilton, James D (1989). ``A New Approach to the Economic Analysis of Nonstationary Time Series and the Business Cycle.'' \textit{Econometrica}, 57(2), 357--384.

\textbf{Core Contribution:}
\begin{itemize}[nosep]
    \item Markov-switching model for regime detection
    \item EM algorithm for parameter estimation
    \item Hidden Markov Model (HMM) framework
\end{itemize}

\textbf{UNTCM Integration:}
\begin{itemize}[nosep]
    \item X4: Markov-Switching UNTCM with regime-specific parameters
    \item E1: Regime Detection algorithm based on HMM
\end{itemize}
\end{tcolorbox}

\begin{tcolorbox}[colback=yellow!5!white,colframe=yellow!75!black,title=Key Paper: Athey \& Imbens (2017)]
\textbf{Citation:} Athey, Susan and Imbens, Guido W (2017). ``The Econometrics of Randomized Experiments.'' \textit{Handbook of Economic Field Experiments}, 1, 73--140.

\textbf{Core Contribution:}
\begin{itemize}[nosep]
    \item ATE and CATE estimation methods
    \item Heterogeneity detection in experiments
    \item Design principles for field experiments
\end{itemize}

\textbf{UNTCM Integration:}
\begin{itemize}[nosep]
    \item E3: Structural estimation from A/B test data
    \item I4: CATE for segment-specific effects
\end{itemize}
\end{tcolorbox}

\begin{tcolorbox}[colback=orange!5!white,colframe=orange!75!black,title=CONTRA Paper: Allcott (2015)]
\textbf{Citation:} Allcott, Hunt (2015). ``Site Selection Bias in Program Evaluation.'' \textit{Quarterly Journal of Economics}, 130(3), 1117--1165.

\textbf{Core Contribution (CONTRA):}
\begin{itemize}[nosep]
    \item Pilot sites overestimate population effects by 2-3x
    \item Selection bias in where programs are tested
    \item Scaling decay is predictable from pilot characteristics
\end{itemize}

\textbf{Implication:} Supports NTM decay calibration; pilot effects should be discounted.
\end{tcolorbox}

\begin{tcolorbox}[colback=orange!5!white,colframe=orange!75!black,title=CONTRA Paper: Vivalt (2020)]
\textbf{Citation:} Vivalt, Eva (2020). ``How Much Can We Generalize from Impact Evaluations?'' \textit{Journal of the European Economic Association}, 18(6), 3045--3089.

\textbf{Core Contribution (CONTRA):}
\begin{itemize}[nosep]
    \item Cross-study variance is high
    \item Prediction accuracy from prior studies is low
    \item Context dependence is substantial
\end{itemize}

\textbf{Implication:} Supports context-modulated decay rates; reinforces need for Ψ-dependent parameters.
\end{tcolorbox}

\subsection{Tipping Points \& Critical Transitions}
\label{subsec:ot-tipping}

Research on early warning signals and critical transitions in complex systems.
\textit{Added 2026-01-20: Supporting UNTCM-E2 (Tipping Point Prediction)}

\begin{tcolorbox}[colback=yellow!5!white,colframe=yellow!75!black,title=Key Paper: Scheffer et al. (2009)]
\textbf{Citation:} Scheffer, Marten et al. (2009). ``Early-Warning Signals for Critical Transitions.'' \textit{Nature}, 461(7260), 53--59.

\textbf{Core Contribution:}
\begin{itemize}[nosep]
    \item Variance increase before tipping (critical slowing down)
    \item Autocorrelation increase as system approaches bifurcation
    \item Flickering between states near regime boundary
\end{itemize}

\textbf{UNTCM Integration:}
\begin{itemize}[nosep]
    \item E2: Early warning indicators for tipping point prediction
    \item Variance and autocorrelation monitoring in $\bar{B}(t)$
\end{itemize}
\end{tcolorbox}

\begin{tcolorbox}[colback=yellow!5!white,colframe=yellow!75!black,title=Key Paper: Scheffer et al. (2012)]
\textbf{Citation:} Scheffer, Marten et al. (2012). ``Anticipating Critical Transitions.'' \textit{Science}, 338(6105), 344--348.

\textbf{Core Contribution:}
\begin{itemize}[nosep]
    \item Generic early warning indicators apply across domains
    \item System-specific indicators needed for precision
    \item Resilience assessment framework
\end{itemize}
\end{tcolorbox}

\subsection{Learning \& Expectations}
\label{subsec:ot-learning}

Research on how agents form and update expectations over time.
\textit{Added 2026-01-20: Supporting UNTCM-X3 (Adaptive Expectations)}

\begin{tcolorbox}[colback=yellow!5!white,colframe=yellow!75!black,title=Key Paper: Muth (1961)]
\textbf{Citation:} Muth, John F (1961). ``Rational Expectations and the Theory of Price Movements.'' \textit{Econometrica}, 29(3), 315--335.

\textbf{Core Contribution:}
\begin{itemize}[nosep]
    \item Expectations are model-consistent on average
    \item Foundation for rational expectations revolution
    \item Learning dynamics in equilibrium
\end{itemize}

\textbf{UNTCM Integration:}
\begin{itemize}[nosep]
    \item X3: $\theta^e$ formation based on model consistency
    \item Adaptive learning with rate $\eta$
\end{itemize}
\end{tcolorbox}

\begin{tcolorbox}[colback=yellow!5!white,colframe=yellow!75!black,title=Key Paper: Sargent (1993)]
\textbf{Citation:} Sargent, Thomas J (1993). ``Bounded Rationality in Macroeconomics.'' Oxford University Press.

\textbf{Core Contribution:}
\begin{itemize}[nosep]
    \item Bounded rationality compatible with learning
    \item Convergence depends on model structure
    \item Heterogeneous learning rates across agents
\end{itemize}

\textbf{UNTCM Integration:}
\begin{itemize}[nosep]
    \item X3: Segment-specific $\eta_s$ (learning rates)
    \item Autonomy-seeking: low $\eta$, Social-oriented: high $\eta$
\end{itemize}
\end{tcolorbox}

\subsection{Behavioral Science Methodology}
\label{subsec:ot-methodology}

Meta-research on behavioral science methods and limitations.
\textit{Added 2026-01-20: Supporting UNTCM Evidence Base (CONTRA)}

\begin{tcolorbox}[colback=orange!5!white,colframe=orange!75!black,title=CONTRA Paper: Bryan et al. (2021)]
\textbf{Citation:} Bryan, Christopher J, Tipton, Elizabeth, and Yeager, David S (2021). ``Behavioural Science Is Unlikely to Change the World Without a Heterogeneity Revolution.'' \textit{Nature Human Behaviour}, 5(8), 980--989.

\textbf{Core Contribution (CONTRA):}
\begin{itemize}[nosep]
    \item Average treatment effects hide crucial heterogeneity
    \item Heterogeneity is systematically underestimated
    \item Need to report and explain effect variation
\end{itemize}

\textbf{UNTCM Implication:}
\begin{itemize}[nosep]
    \item Strongly supports I4 (Segment Heterogeneity) extension
    \item Average UNTCM parameters insufficient; need $\lambda_s, \mu_s, \kappa_s$
    \item Validates segment-specific calibration requirement
\end{itemize}
\end{tcolorbox}

\begin{tcolorbox}[colback=yellow!5!white,colframe=yellow!75!black,title=Key Paper: Gneezy et al. (2011)]
\textbf{Citation:} Gneezy, Uri, Meier, Stephan, and Rey-Biel, Pedro (2011). ``When and Why Incentives (Don't) Work to Modify Behavior.'' \textit{Journal of Economic Perspectives}, 25(4), 191--210.

\textbf{Core Contribution:}
\begin{itemize}[nosep]
    \item Financial incentives can crowd out intrinsic motivation
    \item Crowding-out coefficient $\gamma \approx -0.2$ to $-0.4$
    \item Context determines whether incentives help or hurt
\end{itemize}

\textbf{UNTCM Integration:}
\begin{itemize}[nosep]
    \item I2: T6+T7 crowding interaction ($\gamma_{67} < 0$)
    \item Portfolio design: Avoid simultaneous T6 and T7
\end{itemize}
\end{tcolorbox}

% ============================================================================
% INTEGRATION GUIDELINES
% ============================================================================

%%%%%%%%%%%%%%%%%%%%%%%%%%%%%%%%%%%%%%%%%%%%%%%%%%%%%%%%%%%%%%%%%%%%%%%%%%%%%%%
\section{Key Results}
\label{sec:results}
%%%%%%%%%%%%%%%%%%%%%%%%%%%%%%%%%%%%%%%%%%%%%%%%%%%%%%%%%%%%%%%%%%%%%%%%%%%%%%%

The analysis yields the following key results:

\begin{enumerate}
    \item \textbf{Result 1:} [TODO: Describe main finding]
    \item \textbf{Result 2:} [TODO: Describe secondary finding]
    \item \textbf{Result 3:} [TODO: Describe implication]
\end{enumerate}

\section{Integration Guidelines}
\label{sec:ot-guidelines}

\subsection{Adding Papers}

To add a paper to LIT-OTHER:
\begin{enumerate}
    \item Verify no dedicated author appendix exists
    \item Add to \texttt{data/paper-sources.yaml} with \texttt{ebf\_category: LIT-OTHER}
    \item Document in appropriate thematic cluster above
    \item Note connections to existing EBF appendices
\end{enumerate}

\subsection{Migration Criteria}

Consider creating a dedicated LIT appendix when:
\begin{itemize}
    \item 10+ papers from single author accumulate
    \item Author has Nobel Prize or equivalent recognition
    \item Coherent research program with distinct axioms emerges
    \item Significant EBF-relevant theoretical contribution
\end{itemize}

% ============================================================================
% SUMMARY
% ============================================================================
\section{Summary}
\label{sec:ot-summary}

\begin{tcolorbox}[colback=blue!5!white,colframe=blue!75!black,title=Key Takeaways]
\begin{enumerate}
    \item \textbf{Purpose:} Catch-all for EBF-relevant papers without dedicated appendix
    \item \textbf{Organization:} Thematic clusters rather than author-based
    \item \textbf{Growth:} Papers migrate to dedicated appendices as corpora develop
    \item \textbf{Flexibility:} Allows rapid integration of emerging research
\end{enumerate}

\textbf{Current Holdings:}
\begin{itemize}
    \item Market Morality: 1 paper (Elias et al., 2026)
    \item Choice Architecture: 2 papers (Iyengar \& Lepper 2000, Schwartz 2004)
    \item Cognitive Psychology: 1 paper (Miller 1956)
    \item Prosocial Motivation: 1 paper (Grant 2008)
    \item Organizational Psychology: 1 paper (Meyer \& Allen 1991)
    \item Environmental Behavior: 1 paper (Giner et al. 2025)
    \item Policy Scaling: 1 paper (Soman \& Hossain 2020)
    \item Total: 8 papers
\end{itemize}
\end{tcolorbox}

% ============================================================================
% GLOSSARY
% ============================================================================
\section{Glossary}
\label{sec:ot-glossary}

Key terms defined in Appendix G (REF-GLOSSARY):

\begin{itemize}
    \item \textbf{Repugnant market:} Market where transactions are considered morally objectionable by significant portion of society
    \item \textbf{Contested market:} Market where legitimacy of exchange is debated on ethical grounds
    \item \textbf{Price gouging:} Raising prices significantly above normal levels during emergencies
\end{itemize}

% ============================================================================
% REFERENCES SECTION
% ============================================================================
\section{References}
\label{sec:ot-references}

\nocite{bcm_master}

This appendix collects miscellaneous papers relevant to EBF. Full references are maintained in the master bibliography (\texttt{bibliography/bcm\_master.bib}).

\textbf{Current Papers:}
\begin{itemize}
    \item Elias, J. J., Lacetera, N., \& Macis, M. (2026). The Morality of Market Exchanges. \textit{NBER WP 34647}.
    \item Giner, C., Nauges, C., \& Hassett, K. (2025). Patterns in sustainable food choices. \textit{Food Policy}, 128, 102748.
    \item Iyengar, S. S., \& Lepper, M. R. (2000). When Choice is Demotivating. \textit{JPSP}, 79(6), 995--1006.
    \item Schwartz, B. (2004). \textit{The Paradox of Choice}. Harper Collins.
    \item Miller, G. A. (1956). The Magical Number Seven. \textit{Psychological Review}, 63(2), 81--97.
    \item Grant, A. M. (2008). Prosocial Motivation. \textit{J. Applied Psychology}, 93(1), 48--58.
    \item Meyer, J. P., \& Allen, N. J. (1991). Organizational Commitment. \textit{HRM Review}, 1(1), 61--89.
\end{itemize}

See Appendix G (REF-GLOSSARY) for term definitions.

\end{document}
